% ---------------------------------------------------------------
% Preámbulo: paquetes, estilos de código y macros para los diagramas UML (TikZ)
% ---------------------------------------------------------------
\usepackage[T1]{fontenc}
\usepackage[utf8]{inputenc}
\usepackage[spanish,es-noshorthands,es-tabla,es-nodecimaldot]{babel}
\usepackage{cite}
\usepackage{amsmath,amssymb}
\usepackage{graphicx}
\usepackage[table]{xcolor}
\usepackage{listings}
\usepackage{booktabs}
\usepackage{tabularx}
\usepackage{array}
\usepackage{tikz}
\usetikzlibrary{arrows.meta,positioning,shapes.geometric,shapes.multipart,fit,calc,backgrounds,babel}
\usepackage{url}
\usepackage[hidelinks]{hyperref}

\renewcommand{\IEEEkeywordsname}{Palabras clave}
\renewcommand{\lstlistingname}{Listado}
\def\UrlBreaks{\do\/\do-\do.}

% ---------- Código fuente ----------
\definecolor{fondoCodigo}{RGB}{247,248,250}
\definecolor{clave}{RGB}{0,92,170}
\definecolor{cadena}{RGB}{163,21,21}
\definecolor{comentario}{RGB}{0,128,0}
\definecolor{tipo}{RGB}{38,127,153}

\lstdefinelanguage{TypeScript}{
  keywords={abstract,as,async,await,break,case,catch,class,const,constructor,continue,declare,default,do,else,enum,
    export,extends,false,finally,for,from,function,get,if,implements,import,in,instanceof,interface,keyof,let,
    namespace,new,null,of,override,private,protected,public,readonly,return,satisfies,set,static,super,switch,this,
    throw,true,try,type,typeof,undefined,void,while},
  morekeywords=[2]{string,number,boolean,unknown,never,any,Promise,Record,Partial,Omit,Pick,Readonly,ReadonlyArray,Map},
  sensitive=true,
  morecomment=[l]{//},
  morecomment=[s]{/*}{*/},
  morestring=[b]',
  morestring=[b]",
  morestring=[b]`,
}
\lstdefinelanguage{JavaScript}{
  keywords={async,await,break,case,catch,class,const,constructor,continue,default,do,else,export,extends,false,
    finally,for,from,function,get,if,import,in,instanceof,let,new,null,of,return,set,static,super,switch,this,throw,
    true,try,typeof,undefined,void,while},
  sensitive=true,
  morecomment=[l]{//},
  morecomment=[s]{/*}{*/},
  morestring=[b]',
  morestring=[b]",
  morestring=[b]`,
}
\lstset{
  basicstyle=\ttfamily\scriptsize,
  keywordstyle=\color{clave}\bfseries,
  keywordstyle=[2]\color{tipo},
  stringstyle=\color{cadena},
  commentstyle=\color{comentario}\itshape,
  backgroundcolor=\color{fondoCodigo},
  frame=single,
  rulecolor=\color{gray!40},
  framesep=2pt,
  xleftmargin=3pt,
  xrightmargin=3pt,
  breaklines=true,
  breakatwhitespace=true,
  columns=fullflexible,
  keepspaces=true,
  showstringspaces=false,
  tabsize=2,
  captionpos=b,
  inputencoding=utf8,
  extendedchars=true,
  literate={á}{{\'a}}1 {é}{{\'e}}1 {í}{{\'i}}1 {ó}{{\'o}}1 {ú}{{\'u}}1 {ñ}{{\~n}}1
           {Á}{{\'A}}1 {É}{{\'E}}1 {Í}{{\'I}}1 {Ó}{{\'O}}1 {Ú}{{\'U}}1 {Ñ}{{\~N}}1
           {¿}{{?`}}1 {¡}{{!`}}1,
}

% ---------- Macros UML (TikZ) ----------
\newcommand{\estereotipo}[1]{\guillemotleft #1\guillemotright}

% Actor (monigote). #1 nombre, #2 posición "x,y", #3 etiqueta.
% Define las coordenadas #1-e (este), #1-o (oeste) y el nodo #1-etq (etiqueta).
\newcommand{\actorUML}[3]{%
  \begin{scope}[shift={(#2)}]
    \draw[thick] (0,0.62) circle (0.17);
    \draw[thick] (0,0.45) -- (0,-0.05);
    \draw[thick] (-0.32,0.28) -- (0.32,0.28);
    \draw[thick] (0,-0.05) -- (-0.24,-0.5);
    \draw[thick] (0,-0.05) -- (0.24,-0.5);
    \node[below, font=\small, inner sep=2pt] (#1-etq) at (0,-0.52) {#3};
    \coordinate (#1-e) at (0.42,0.2);
    \coordinate (#1-o) at (-0.42,0.2);
  \end{scope}}

% Icono de componente UML en la esquina superior derecha de un nodo
\newcommand{\iconoComponente}[1]{%
  \draw[fill=white] ($(#1.north east)+(-0.40,-0.10)$) rectangle ++(0.26,-0.30);
  \draw[fill=white] ($(#1.north east)+(-0.46,-0.15)$) rectangle ++(0.12,-0.06);
  \draw[fill=white] ($(#1.north east)+(-0.46,-0.27)$) rectangle ++(0.12,-0.06);}

% Estilos de los diagramas de clases
\tikzset{
  clases/.style={
    clase/.style={draw, thick, fill=white, anchor=north, rectangle split, rectangle split parts=3,
      rectangle split part fill={teal!14,white,white}, text width=3.4cm, align=left,
      font=\scriptsize, inner sep=2.5pt, execute at begin node={\setlength{\parindent}{0pt}}},
    interfaz/.style={clase, rectangle split parts=2, rectangle split part fill={orange!14,white}},
    simple/.style={clase, rectangle split parts=2, rectangle split part fill={teal!14,white}},
    herencia/.style={thick, -{Triangle[open, length=2.6mm, width=2.6mm]}},
    realiza/.style={dashed, -{Triangle[open, length=2.6mm, width=2.6mm]}},
    asociacion/.style={-{Stealth[length=2mm]}},
    depende/.style={dashed, -{Stealth[length=2mm]}},
    composicion/.style={{Diamond[length=3mm, width=2mm]}-},
    agregacion/.style={{Diamond[open, length=3mm, width=2mm]}-},
    rol/.style={font=\tiny, inner sep=1pt},
  },
}
% \clase{nombre}{x,y}{encabezado}{atributos}{métodos}
\newcommand{\clase}[5]{\node[clase] (#1) at (#2)
  {{\centering #3\par}\nodepart{two}{\raggedright #4\par}\nodepart{three}{\raggedright #5\par}};}
% \interfaz / \simple {nombre}{x,y}{encabezado}{cuerpo}
\newcommand{\interfaz}[4]{\node[interfaz] (#1) at (#2) {{\centering #3\par}\nodepart{two}{\raggedright #4\par}};}
\newcommand{\simple}[4]{\node[simple] (#1) at (#2) {{\centering #3\par}\nodepart{two}{\raggedright #4\par}};}
